\documentclass[11pt]{article}
\usepackage[landscape, margin=0.5in]{geometry}
\usepackage{booktabs}
\usepackage{longtable}
\usepackage{array}
\usepackage{numprint}
\usepackage{float}
\usepackage{xcolor}
\usepackage{amsmath}
\usepackage{titlesec}

% Custom column types for better formatting
\newcolumntype{R}[1]{>{\raggedleft\arraybackslash}m{#1}}
\newcolumntype{L}[1]{>{\raggedright\arraybackslash}m{#1}}
\newcolumntype{C}[1]{>{\centering\arraybackslash}m{#1}}

% Set up number formatting
\nprounddigits{0}

\title{Analysis of BEA Code Mapping Differences: \\
Original vs. Alternative Census Mappings}
\author{Federal Reserve Bank of Atlanta}
\date{\today}

\begin{document}

\maketitle

\section{Executive Summary}

The comparison between original Schott-based HS-to-BEA mappings and the new Alternative Census Bureau mappings reveals substantial differences in import value coverage. The alternative mapping captures \textbf{\$148.2 billion} less in import values than the original mapping, representing a \textbf{4.54\%} reduction in total coverage. This analysis examines the specific BEA codes, countries, and industries driving these differences.

\section{Overall Impact Summary}

\begin{table}[H]
\centering
\caption{Total Import Value Comparison Across BEA Aggregation Levels}
\begin{tabular}{lrrr}
\toprule
\textbf{Aggregation Level} & \textbf{Original (\$B)} & \textbf{Alternative (\$B)} & \textbf{Difference (\$B)} \\
\midrule
Detail & 3,267.4 & 3,119.2 & -148.2 \\
U.Summary & 3,267.4 & 3,119.2 & -148.2 \\
Summary & 3,267.4 & 3,119.2 & -148.2 \\
Sector & 3,267.4 & 3,119.2 & -148.2 \\
\bottomrule
\end{tabular}
\label{tab:overall_impact}
\end{table}

The consistent \$148.2 billion difference across all aggregation levels indicates that the mapping differences occur at the most granular (detail) level and propagate upward through the BEA hierarchy.

\section{Primary Sources of Differences}

\subsection{BEA Code S00300: "Noncomparable Imports"}

The overwhelming majority of differences (>99\%) stem from the complete absence of BEA code S00300 in the alternative Census mapping. This code represents "noncomparable imports" - trade flows that cannot be matched to specific NAICS industries.

\begin{table}[H]
\centering
\caption{Top 15 Countries with S00300 Losses}
\begin{tabular}{lR{2.5cm}R{2.5cm}}
\toprule
\textbf{Country} & \textbf{Original Value (\$M)} & \textbf{Lost Value (\$M)} \\
\midrule
Canada & 24,918.2 & 24,918.2 \\
Mexico & 13,990.4 & 13,990.4 \\
Germany & 12,440.1 & 12,440.1 \\
China & 11,303.1 & 11,303.1 \\
United Kingdom & 8,024.2 & 8,024.2 \\
Ireland & 7,460.5 & 7,460.5 \\
France & 6,661.5 & 6,661.5 \\
Italy & 5,426.7 & 5,426.7 \\
Netherlands & 5,044.1 & 5,044.1 \\
Japan & 4,844.5 & 4,844.5 \\
Singapore & 4,808.4 & 4,808.4 \\
Switzerland & 3,759.5 & 3,759.5 \\
Brazil & 3,032.2 & 3,032.2 \\
Belgium & 2,965.9 & 2,965.9 \\
Hong Kong & 2,727.0 & 2,727.0 \\
\midrule
\textbf{Total (15 countries)} & \textbf{116,405.8} & \textbf{116,405.8} \\
\bottomrule
\end{tabular}
\label{tab:s00300_losses}
\end{table}

The S00300 category affects \textbf{231 countries} and represents imports with HS codes typically in the 910000, 930000, 980000, and 990000 ranges - special trade categories that include diplomatic goods, returned goods, and other non-standard trade flows.

\subsection{Battery Manufacturing Reclassification (335912 $\leftrightarrow$ 335911)}

A systematic reclassification appears to occur between "Primary Battery Manufacturing" (335912) and "Storage Battery Manufacturing" (335911), with the alternative mapping favoring the latter classification.

\begin{table}[H]
\centering
\caption{Battery Manufacturing Code Transfers (Top Countries)}
\begin{tabular}{lR{2cm}R{2cm}R{2cm}}
\toprule
\textbf{Country} & \textbf{335912 Loss (\$M)} & \textbf{335911 Gain (\$M)} & \textbf{Net Transfer (\$M)} \\
\midrule
China & 378.0 & 378.0 & 0.0 \\
Indonesia & 130.8 & 130.8 & 0.0 \\
Singapore & 125.0 & 125.0 & 0.0 \\
Japan & 95.5 & 95.5 & 0.0 \\
Belgium & 77.3 & 77.3 & 0.0 \\
Israel & 57.3 & 57.3 & 0.0 \\
Canada & 48.0 & 48.0 & 0.0 \\
Vietnam & 44.3 & 44.3 & 0.0 \\
\midrule
\textbf{Total} & \textbf{956.2} & \textbf{956.2} & \textbf{0.0} \\
\bottomrule
\end{tabular}
\label{tab:battery_reclassification}
\end{table}

This reclassification suggests the alternative Census mapping may use updated or different criteria for distinguishing between primary and storage battery manufacturing, but results in zero net impact on total import values.

\subsection{Automobile Manufacturing Losses (336111)}

Some automobile imports are completely lost in the alternative mapping, affecting major automotive trading partners:

\begin{table}[H]
\centering
\caption{Automobile Manufacturing (336111) Losses}
\begin{tabular}{lR{2.5cm}}
\toprule
\textbf{Country} & \textbf{Lost Value (\$M)} \\
\midrule
Korea, South & 717.8 \\
Other countries & 69.2 \\
\midrule
\textbf{Total} & \textbf{787.0} \\
\bottomrule
\end{tabular}
\label{tab:auto_losses}
\end{table}

\section{Record Count Analysis}

The differences in record counts reveal the scope of mapping coverage changes:

\begin{table}[H]
\centering
\caption{Record Count Comparison by Aggregation Level}
\begin{tabular}{lR{1.5cm}R{1.5cm}R{1.5cm}R{1.5cm}}
\toprule
\textbf{Level} & \textbf{Both} & \textbf{Alt-Only} & \textbf{Orig-Only} & \textbf{Total} \\
\midrule
Detail & 22,178 & 6 & 260 & 22,444 \\
U.Summary & 5,930 & 0 & 236 & 6,166 \\
Summary & 3,905 & 0 & 231 & 4,136 \\
Sector & 773 & 0 & 231 & 1,004 \\
\bottomrule
\end{tabular}
\label{tab:record_counts}
\end{table}

The pattern shows that:
\begin{itemize}
\item The alternative mapping covers most existing country-industry combinations (>95\% overlap)
\item Very few new combinations are introduced (only 6 at detail level)
\item 231-260 original combinations are lost across aggregation levels
\end{itemize}

\section{Industry Classification Implications}

\subsection{Noncomparable Imports (S00300)}

The complete absence of S00300 classifications in the alternative mapping represents the most significant limitation. These imports typically include:
\begin{itemize}
\item HS Code 910000: Special classifications
\item HS Code 930000: Special transactions  
\item HS Code 980000: Estimated imports
\item HS Code 990000: Returned goods
\end{itemize}

The inability to map these special categories suggests the Alternative Census mapping may be more restrictive in handling non-standard trade flows.

\subsection{Manufacturing Sector Refinements}

The battery manufacturing reclassification (335912 $\rightarrow$ 335911) and automobile manufacturing losses (336111) indicate that the alternative mapping:
\begin{itemize}
\item May use more recent or refined industry classification criteria
\item Could reflect changes in manufacturing processes or product definitions
\item Might have different interpretations of HS code-to-NAICS relationships
\end{itemize}

\section{Geographic Distribution of Impact}

\begin{table}[H]
\centering
\caption{Countries Most Affected by Mapping Differences (Top 20)}
\begin{tabular}{lR{2.5cm}R{2cm}}
\toprule
\textbf{Country} & \textbf{Total Absolute Difference (\$M)} & \textbf{\% of Total Loss} \\
\midrule
Canada & 24,918.2 & 16.8\% \\
Mexico & 13,990.4 & 9.4\% \\
Germany & 12,440.1 & 8.4\% \\
China & 11,303.1 & 7.6\% \\
United Kingdom & 8,024.2 & 5.4\% \\
Ireland & 7,460.5 & 5.0\% \\
France & 6,661.5 & 4.5\% \\
Italy & 5,426.7 & 3.7\% \\
Netherlands & 5,044.1 & 3.4\% \\
Japan & 4,844.5 & 3.3\% \\
Singapore & 4,808.4 & 3.2\% \\
Switzerland & 3,759.5 & 2.5\% \\
Brazil & 3,032.2 & 2.0\% \\
Belgium & 2,965.9 & 2.0\% \\
Hong Kong & 2,727.0 & 1.8\% \\
Taiwan & 2,591.0 & 1.7\% \\
Korea, South & 2,265.5 & 1.5\% \\
Vietnam & 1,942.9 & 1.3\% \\
Malaysia & 1,942.8 & 1.3\% \\
India & 1,479.1 & 1.0\% \\
\midrule
\textbf{Top 20 Total} & \textbf{126,627.4} & \textbf{85.4\%} \\
\bottomrule
\end{tabular}
\label{tab:geographic_impact}
\end{table}

The losses are concentrated among major U.S. trading partners, with the top 20 countries accounting for 85.4\% of total differences.

\section{Conclusions and Recommendations}

\subsection{Key Findings}

\begin{enumerate}
\item The Alternative Census mapping reduces import coverage by \$148.2 billion (4.54\%) compared to the original Schott-based mapping.

\item The primary driver is the complete absence of "noncomparable imports" (S00300) classification, affecting 231 countries.

\item Battery manufacturing shows systematic reclassification between primary (335912) and storage (335911) categories with zero net impact.

\item Some automobile manufacturing imports (336111) are lost, particularly from South Korea.
\end{enumerate}

\subsection{Implications for Research}

\begin{itemize}
\item \textbf{Trade Coverage}: The 4.54\% reduction in import coverage may affect comprehensive trade analyses.

\item \textbf{Special Categories}: Research requiring complete coverage of all import types may be limited by the absence of noncomparable imports.

\item \textbf{Industry Analysis}: Manufacturing sector analyses may benefit from the apparent refinements in battery and automotive classifications.

\item \textbf{Country Studies}: Analyses of major trading partners (especially Canada, Mexico, Germany, China) will show notable differences.
\end{itemize}

\subsection{Recommendations}

\begin{enumerate}
\item \textbf{Assess Research Requirements}: Determine whether the 4.54\% coverage reduction and loss of special trade categories is acceptable for the intended research objectives.

\item \textbf{Investigate S00300 Handling}: Explore whether noncomparable imports can be incorporated into the alternative mapping or if their exclusion is by design.

\item \textbf{Validate Industry Reclassifications}: Confirm that battery and automotive sector changes reflect improved accuracy rather than mapping errors.

\item \textbf{Document Limitations}: Clearly communicate the coverage limitations when presenting results based on the alternative mapping.

\item \textbf{Consider Hybrid Approach}: Evaluate whether combining aspects of both mappings could optimize coverage while maintaining accuracy.
\end{enumerate}

The Alternative Census mapping represents a more restrictive but potentially more accurate approach to HS-to-BEA code mapping, with trade-offs in coverage that must be weighed against research objectives.

\end{document}